\section{Conclusion}
\label{sec:conclusion}

Constructive and destructive regimes connect the interaction of successive gradients to curvature, accumulation, and changes in the learned function. Their interpretation depends on the object and timescale being measured. Negative adjacent interaction can coexist with positive total cross-step interaction and lower loss. Under smooth GD, two existing gradients reveal an exact step-averaged directional Hessian scalar, with explicit sampling and finite-step qualifications. Under Adam, the same observable describes a mixed gradient response; exhaustive timing checks show useful structure without establishing a universal next-step alarm.

Fixed-probe tracing makes the functional expression of these regimes visible. Across the three controlled cases, representation motion concentrates in leading modes, shares temporal structure across measured depth, and repeatedly reverses along fitted output patterns. A local affine-displacement cap then tests the usefulness of that interpretation. Every tested cap quantile reduces mean training loss and positive variation relative to its matched unchanged-Adam control, while the preselected median cap retains frequent reversals. Its ResNet endpoint degradation, held-out losses, and implementation cost prevent interpreting this result as universal optimizer superiority or improved generalization.

The practical implication is to distinguish reversal from excessive movement and to evaluate intervention effects jointly with learning progress. The next evidential steps are matched independent initializations, prospective diagnostics with fixed horizons and loss-history baselines, independent probes and quiet-window controls, and representation-aware comparators with matched tuning budgets. Stochastic deployment and larger-scale evaluation remain separate tests. Within the present scope, the framework provides a reproducible route from observable gradient interaction to a bounded, experimentally evaluated functional control.
